% ---
% id: protocol-compiler/ex:two-party-shape
% kind: ex
% title: "Declaration of a two-party protocol with an untrusted source"
% summary: >-
%   An untrusted source serves two measuring parties; the measured system is a
%   product of two factors with their own sector maps and families, and no
%   product matrix is formed. A party-local entry bounds Bob's flag weight by
%   1/8 conditionally on every state of Alice's measured factor.
% prominence: standard
% status: accepted
% uses: []
% mentions: [def:declaration, def:target, def:feasible, def:homog, def:exact, def:spec, transmitter/rem:da1, rem:party-local-soundness]
% assumes: []
% refs: []
% provenance: newly-derived
% lean:
% notes: "The numbers are chosen dyadic so that every exact predicate of def:exact holds of them as written."
% ---

The example writes out the general format of
Definition~\ref{def:declaration} for a protocol in which an untrusted source
serves two \tref{measuring parties}. It illustrates the format and asserts nothing
about any physical protocol.

\begin{example}\label{ex:two-party-shape}
\emph{Systems and sector map.}
\tref[untrusted-source mode]{Untrusted-source mode}, $d_A=1$: the
\tref{marginal slot} holds $\nomarg$, $\Xset=\{x_0\}$ and $M_{x_0}=[1]$
(\xref{transmitter/rem:da1}). The sector map has no factor tagged $A$ and two
factors tagged $\Qsys$, in this order: $f=\mathsf a$ with party label
$\mathsf A$ and $f=\mathsf b$ with party label $\mathsf B$
(Definition~\ref{def:declaration}\,(DX1)). Each factor has a \tref{data sector} of
dimension $2$ and a flag sector of dimension $1$,
$\SecSet^{f}=(\mathrm{d},\mathrm{fl})$, $d_{f,\mathrm d}=2$,
$d_{f,\mathrm{fl}}=1$. The product map has the four sectors
$(\mathrm d,\mathrm d),(\mathrm d,\mathrm{fl}),(\mathrm{fl},\mathrm d),
(\mathrm{fl},\mathrm{fl})$ in that order, the last factor fastest, of
dimensions $4,2,2,1$; $|\mathbb{T}|=1$, so these are also the
\tref[joint sector]{joint sectors}, and a state $\omega$ is a list of four blocks of those sizes.

\emph{Family.} Each factor declares its own family over
$\Vset^{f}=(0,1,\mathrm{fl})$, in this order:
$F^{f,(\mathrm d)}_0=\mathrm{diag}(1,0)$,
$F^{f,(\mathrm d)}_1=\mathrm{diag}(0,1)$, $F^{f,(\mathrm d)}_{\mathrm{fl}}=0$,
and $F^{f,(\mathrm{fl})}_0=F^{f,(\mathrm{fl})}_1=[0]$,
$F^{f,(\mathrm{fl})}_{\mathrm{fl}}=[1]$. Each is \tref{exactly admissible}
for its factor's map, its two completeness sums rounding nowhere
(Definition~\ref{def:exact}, per measured factor). The \tref{outcome alphabet} is
$\Vset=\Vset^{\mathsf a}\times\Vset^{\mathsf b}$ with nine letters, and
$F_{(v_{\mathsf a},v_{\mathsf b})}=F^{\mathsf a}_{v_{\mathsf a}}\otimes
F^{\mathsf b}_{v_{\mathsf b}}$ by the rule \eqref{eq:productfamily}; its block
on $(\mathrm d,\mathrm{fl})$, for instance, is
$F^{\mathsf a,(\mathrm d)}_{v_{\mathsf a}}\otimes
F^{\mathsf b,(\mathrm{fl})}_{v_{\mathsf b}}$, a $2\times2$ matrix that is
written down by nobody. Nine $9\times9$ arrays are never formed.

\emph{\tref[party partition]{Party partition} and announcements.}
$\rd{\mathsf A}(x_0,v)=v_{\mathsf a}$ and $\rd{\mathsf B}(x_0,v)=v_{\mathsf b}$,
which is \emph{not} the \tref{default party partition}: under the default one
both parties' data would sit in $\rd{\mathsf B}$ and party computability would
be vacuous. A transcript with $L=2$: $\pi=(\mathsf A,\mathsf B)$, each party
announcing whether its outcome is a flag,
$a_1(v_{\mathsf a})=a_2(v_{\mathsf b},\kappa_1)=\mathrm{fl}$ on the flag
letter and $\mathrm{ok}$ otherwise. Sifting passes on
$\kappa=(\mathrm{ok},\mathrm{ok})$; $\hold=\mathsf A$, $d_Z=2$, and the key
rule $r^{0}_{\mathsf A}$ of the specification's public part
(Definition~\ref{def:spec}\,(d)), into $\Zval=\{0,1\}$, is
$r^{0}_{\mathsf A}(v_{\mathsf a},\kappa)=v_{\mathsf a}$ for
$v_{\mathsf a}\in\{0,1\}$ and $r^{0}_{\mathsf A}(\mathrm{fl},\kappa)=0$, a
value no passing round reads, since $v_{\mathsf a}=\mathrm{fl}$ announces
$\kappa_1=\mathrm{fl}$. Party computability, \eqref{eq:partycomp}, requires
the raw key to be a function of Alice's own readout and the transcript, as it
is.

\emph{A party-local entry.} Let $\hat O$ be the block list on factor
$\mathsf b$ with $\hat O^{(\mathrm d)}=0$ and $\hat O^{(\mathrm{fl})}=[1]$,
the projector onto Bob's flag sector, and declare
$(\hat O,\preceq,\tfrac18;\mathsf b)$ with the tag party-local. By the last
line of \eqref{eq:feas-general}, with $s_{f^c}=s_{\mathsf a}$, it reads, for
each of Alice's two sectors $s_{\mathsf a}$,
\begin{equation*}
  \tr_{\Qsys_{\mathsf b,\mathrm{fl}}}\omega^{(s_{\mathsf a},\mathrm{fl})}
  -\tfrac18\Big(\tr_{\Qsys_{\mathsf b,\mathrm d}}
  \omega^{(s_{\mathsf a},\mathrm d)}
  +\tr_{\Qsys_{\mathsf b,\mathrm{fl}}}\omega^{(s_{\mathsf a},\mathrm{fl})}\Big)
  \;\preceq\;0
  \qquad\text{on }\Hs{\Qsys_{\mathsf a,s_{\mathsf a}}} ,
\end{equation*}
a $2\times2$ and a $1\times1$ condition: Bob's flag weight is at most
$\tfrac18$ \emph{conditionally on every state of Alice's measured factor
$\mathsf a$}. The same
entry without the tag reads
$\tr[(\idon{}\otimes\hat O)\omega]\le\tfrac18$, one real inequality, which is
the trace of the sum of the two conditions above and does not imply them: it
allows all of Bob's flag weight to sit where Alice's measured factor is in
one fixed state. Which reading the \tref{producer} of the entries may emit is decided by
its certificate (Definition~\ref{def:homog},
Remark~\ref{rem:party-local-soundness}), not here.
\end{example}
